% =========================
% Industry CV (UK/Ireland) - Single-page friendly LaTeX template
% Author: You (customised by Dr Susam Boral)
% =========================

\documentclass[10pt,a4paper]{article}

% ---------- Page + Fonts ----------
\usepackage[a4paper,margin=1in]{geometry}
\usepackage[T1]{fontenc}
\usepackage[utf8]{inputenc}
\usepackage{lmodern} % ATS-friendly font
\usepackage{microtype}

% ---------- Utilities ----------
\usepackage{enumitem}
\usepackage{hyperref}
\usepackage{xcolor}
\usepackage{titlesec}

\hypersetup{
  colorlinks=true,
  urlcolor=blue,
  linkcolor=black
}

% ---------- Section styling (simple, ATS-friendly) ----------
\titleformat{\section}{\large\bfseries}{}{0em}{}[\titlerule]
\titlespacing*{\section}{0pt}{8pt}{6pt}

% ---------- Lists ----------
\setlist[itemize]{leftmargin=1.2em, itemsep=2pt, topsep=2pt}
\setlist[enumerate]{leftmargin=1.2em, itemsep=2pt, topsep=2pt}

% ---------- Custom helpers ----------
\newcommand{\cvline}[2]{\textbf{#1} \hfill #2}
\newcommand{\role}[3]{\textbf{#1} \hfill \textit{#2}\\ \textit{#3}}
\newcommand{\tightsep}{\vspace{2pt}}

\pagenumbering{gobble}

\begin{document}

% =========================
% Header
% =========================
{\LARGE \textbf{Dr Susam Boral}}\\[-2pt]
Dublin, Ireland \;|\; +353 87 031 0481 \;|\; \href{mailto:susamboral@outlook.com}{susamboral@outlook.com} \;|\;
\href{https://scholar.google.com/citations?user=lxZzGV8AAAAJ&hl=en}{Google Scholar}
\vspace{6pt}

% =========================
% Professional Summary (3–4 lines max)
% =========================
\section*{Professional Summary}
Research Fellow with a PhD and postdoctoral experience in applied mathematics and computational modelling, working on industry-relevant R\&D problems. Experienced in developing and validating numerical models, translating complex results into clear technical insight, and collaborating with multidisciplinary teams and external partners. Interested in an industry research role focused on early-stage technology development, collaborative programmes, and turning research outcomes into practical capability.

% =========================
% Key Skills (tailor to each job)
% =========================
\section*{Key Skills}
\begin{itemize}
  \item \textbf{R\&D \& Innovation:} Early-stage research, feasibility studies, technology scouting (literature/patents), translating research into practical capability
  \item \textbf{Modelling \& Analysis:} Numerical modelling, predictive simulation, uncertainty and sensitivity analysis, model verification/validation
  \item \textbf{Programming \& Tools:} Python, MATLAB, scientific computing, data analysis, reproducible workflows and documentation
  \item \textbf{Research Collaboration:} Working with academia/industry partners, proposal support, workshop coordination, technical presentations
  \item \textbf{Communication \& Leadership:} Technical writing, stakeholder communication, mentoring/supervision, cross-functional teamwork
\end{itemize}
% =========================
% Experience (industry-style bullets with outcomes)
% =========================
\section*{Experience}
\role{Research Fellow, Trinity College Dublin}{Nov 2024 -- Present}{Dublin, Ireland}
\begin{itemize}
  \item Developed predictive and simulation-based models to analyse complex systems under uncertainty, supporting risk-aware decision frameworks.
  \item Performed statistical and sensitivity analyses to evaluate model robustness and scenario-based system behaviour.
  \item Communicated analytical insights and modelling results to multidisciplinary stakeholders, supporting data-driven decision making.
\end{itemize}

\role{Postdoctoral Fellow, Harbin Engineering University}{Jun 2023 -- Oct 2024}{Harbin, China}
\begin{itemize}
  \item Conducted advanced modelling studies in hydrodynamics/hydroelasticity; validated results through analytical checks and numerical benchmarking.
  \item Collaborated across institutions; contributed to publications and project planning with clear documentation and versioned code.
\end{itemize}

\role{Research Associate, IIT Kharagpur}{Mar 2023 -- May 2023}{Kharagpur, India}
\begin{itemize}
  \item Supported research delivery through modelling, analysis, and preparation of technical material for dissemination.
\end{itemize}

% =========================
% Selected Projects (very useful for consultancy applications)
% =========================
\section*{Selected Projects}

\textbf{InDyCORE: Dynamic cables for offshore renewable energy (SEAI-funded)} \hfill \textit{Trinity College Dublin}
\begin{itemize}
  \item Contributed to applied R\&D by developing and validating computational models for complex system behaviour under real-world loading.
  \item Used sensitivity/parameter studies to improve understanding of key drivers and support early-stage design and technical readiness.
  \item Worked with multidisciplinary researchers and partners to deliver documented results, presentations, and technical outputs.
\end{itemize}

\textbf{Computational Modelling and AI/Data Science Training} \hfill \textit{2024--2025}
\begin{itemize}
  \item Completed Executive PG Certification in Data Science and AI; strengthened capability in data analysis, statistical modelling, and Python workflows.
  \item Applied modelling and data-driven methods to research problems, focusing on validation, robustness checks, and interpretable results.
\end{itemize}
% =========================
% Education (keep short for industry)
% =========================
\section*{Education}
\cvline{Executive PG Certification, Data Science \& AI}{Jun 2024 -- Jul 2025}\\
iHUB DivyaSampark, IIT Roorkee and Intellipaat

\tightsep
\cvline{PhD, Fluid dynamics}{2022}\\
Indian Institute of Technology Kharagpur

\tightsep
\cvline{MSc, Applied Mathematics}{2018}\\
University of Calcutta

\tightsep
\cvline{BSc, Mathematics (Hons.)}{2016}\\
University of Calcutta

% =========================
% Certifications / Training (optional)
% =========================
\section*{Training \& Professional Development}
\begin{itemize}
  \item Mathematics for Machine Learning (Coursera) \hfill 2021
  \item Short courses/workshops in coastal/offshore engineering and hydroelasticity \hfill 2018--2022
\end{itemize}

% =========================
% Publications (keep short for industry; link out)
% =========================
\section*{Publications (Selected)}
\begin{itemize}
  \item Selected peer-reviewed publications in applied mathematics, hydrodynamics, and hydroelasticity (full list: \href{https://scholar.google.com/citations?user=lxZzGV8AAAAJ&hl=en}{Google Scholar})
\end{itemize}

% =========================
% Awards (short)
% =========================
\section*{Awards}
\begin{itemize}
  \item CSIR Senior Research Fellowship, India \hfill 2020--2022
  \item CSIR Junior Research Fellowship, India \hfill 2018--2020
  \item First rank in MSc (Applied Mathematics), University of Calcutta \hfill 2018
    \item First rank in BSc (Mathematics Honours), University of Calcutta \hfill 2018
\end{itemize}

% =========================
% Leadership / Service (optional; keep concise)
% =========================
\section*{Leadership \& Service}
\begin{itemize}
  \item Organiser, International Workshop on Mathematics of Sea Ice and Ice Sheets (MOSSI) \hfill 2023--2024
  \item Visiting Researcher, Isaac Newton Institute, Cambridge \hfill 2023
\end{itemize}

% =========================
% Languages
% =========================
\section*{Languages}
Bengali (native), English (fluent), Hindi (fluent)

% =========================
% Notes (remove before final)
% =========================
% Tips:
% 1) Keep to 1–2 pages for industry.
% 2) Remove long publication lists; keep a link instead.
% 3) Remove personal details like DOB, gender, marital status, full address.
% 4) Tailor Key Skills + Selected Projects to the job description each time.

\end{document}